\setcounter{neje}{67}
\begin{nejer}
Sean $(X, \mathcal{A}, \mu)$ y $(X, \mathcal{A}', \mu')$ dos espacios de medida con $\mathcal{A} \subseteq \mathcal{A}'$ y $\mu'|_{\mathcal{A}} = \mu$. Pruebe que $\mathcal{L}^1(\mu) \subseteq \mathcal{L}^1(\mu')$ y
$$
\int_X f \dd{\mu} = \int_X f \dd{\mu'}
$$
para cada $f \en \mathcal{L}^1(\mu)$.
% \emph{Sugerencia.} Pruebe lo siguiente:
% \begin{enumerate}[label = (\roman*)]
%     \item Sea $f\colon X \to [0, \infty]$ una función $\mathcal{A}$-medible y defina
%     $$
%     f_\delta(x) := \begin{cases} 
%     0 & \text{si } f(x) \le \delta, \\
%     f(x) & \text{si } \delta < f(x) \le \delta^{-1}, \\
%     \delta^{-1} & \text{si } f(x) > \delta^{-1}.
%     \end{cases}
%     $$
%     Entonces $f_\delta$ es $\mathcal{A}$-medible para todo $\delta > 0$ y
%     $$
%     \lim_{\delta \to 0} \int_X f_\delta \dd{\mu} = \int_X f \dd{\mu}.
%     $$
    

\end{nejer}
\begin{sol}
Este ejercicio venía con una sugerencia que no seguiré.\vspace{1em}\\
\emph{Probémoslo para las funciones escalonadas no negativas:}\\
Sea $s\colon X \to [0,\infty]$ escalonada $\mathcal{A}$-medible, es decir, $s = \sum^m_{k=1}a_k\chi_{A_k}$ con $a_k\en[0,\infty]$ y $A_k\en\mathcal{A}$. Como $\mathcal{A} \subseteq \mathcal{A}'$, entonces cada $A_k \en\mathcal{A}'$, así que $s$ es también $\mathcal{A}'$-medible. Además, $\mu'|_\mathcal{A} = \mu$, así
\[
\int_X s \dd{\mu'} = \sum^m_{k=1} a_k\mu'(A_k) = \sum^m_{k=1} a_k\mu(A_k) = \int_X s \dd{\mu}
\]
\\
\emph{Probémoslo para funciones medibles no negativas}:\\
Si $f\colon X \to [0,\infty]$ es $\mathcal{A}$-medible, entonces existe una sucesión $\set{s_n}_{n\en\mathbb{N}}$ monótona de funciones escalonadas, $\mathcal{A}$-medibles y no negativas, tal que $s_n \to f$. El teorema de convergencia monótona garantiza que 
\[
\lim_{n\to\infty} \int_Xs_n\dd{\mu} = \int_X f \dd{\mu}
\]
No obstante, lo probado al inicio garantiza que cada $s_n$ es $\mathcal{A}'$-medible, así que $f$ también será $\mathcal{A}'$-medible por ser límite puntual de funciones $\mathcal{A}'$-medibles. Más aún, por convergencia monótona:
\[
\lim_{n\to\infty} \int_Xs_n\dd{\mu'} = \int_X f \dd{\mu'}
\]
No obstante, ya vimos que $\int_X s_n \dd{\mu} = \int_X s_n \dd{\mu'}$, por lo cual:
\[
\int_X f \dd{\mu} = \lim_{n\to\infty} \int_Xs_n\dd{\mu} = \lim_{n\to\infty} \int_Xs_n\dd{\mu'} = \int_X f \dd{\mu'}
\]
\\
\emph{Probemos el resultado}:\\
Sea $f\en\mathcal{L}^1(\mu)$. Se sigue que $\abs{f},f^+,f^-$ son $\mathcal{A}$-medibles (no negativas), así que lo recién probado garantiza que también son $\mathcal{A}'$-medibles y además
\[
\int_X \abs{f} \dd{\mu} = \int_X \abs{f} \dd{\mu'},
\quad
\int_X f^+ \dd{\mu} = \int_X f^+ \dd{\mu'}
\quad\text{y}\quad
\int_X f^- \dd{\mu} = \int_X f^- \dd{\mu'}
\]
Por un lado, esto implica que $f = f^+ - f^-$ es $\mathcal{A}'$-medible al ser suma de funciones $\mathcal{A}'$-medibles. Por otro lado, garantiza que $\int_X \abs{f} \dd{\mu'} = \int_X \abs{f} \dd{\mu} < \infty$, por lo que $f\en\mathcal{L}^1(\mu')$. Más aún, sustituyendo lo anterior en la definición de integral real-valuada:
\[
\int_X f \dd{\mu'} 
= \int_X f^+ \dd{\mu'} - \int_X f^- \dd{\mu'}
= \int_X f^+ \dd{\mu} - \int_X f^- \dd{\mu}
= \int_X f \dd{\mu}
\]
Hemos demostrado que $\mathcal{L}^1(\mu) \subseteq \mathcal{L}^1(\mu')$ y $\int_X f \dd{\mu} = \int_X f \dd{\mu'}$.
\end{sol}